% Tutorial, Chapter 4 -- Working with Formulas: Lessons 12 to 15 (AN-17).
% Figures: screens/tl12..tl15.k123 on the data volume of stage 3 (BUDGET
% saved).  Sources (Law 1): os/i123/core/formula.c (absolute references,
% $NAME), os/i123/core/range.c (range names; the control panel shows the
% range a name stands for), os/i123/core/recalc.c (manual recalculation,
% CALC), os/i123/core/eval.c (@IF; strings in quotation marks).
% GATES: test-i123-formula test-i123-recalc-order test-i123-range-menu test-pubs-screens test-pubs-i123
\Gates{test-i123-formula, test-i123-recalc-order, test-i123-range-menu,
  test-pubs-screens, test-pubs-i123}
\chapter{Working with Formulas}

In this chapter you add each supply's share of the budget, name a range,
recalculate by hand, and add a test that marks the expensive quarters.
Begin by retrieving BUDGET: press \kname{SLASH}, type \type{f} and
\type{r}, type \type{BUDGET} and press \kname{ENTER}.

\section{Lesson 12: Absolute References}
\ix{absolute reference}
Each supply's share of the year's spending is its total divided by the
total in F9.  When you copy such a formula, the reference to F9 must not
move.  An \term{absolute reference}, with dollar signs, \texttt{\$F\$9},
always refers to the same cell.

\begin{Keystrokes}
\Do{Press}{\fkey{Goto}{F5}}{and type \type{g4}, then press \kname{ENTER}}
\Do{Type}{\type{"Share}}{and press \kname{DOWN}}
\Do{Type}{\type{+f5/\$f\$9}}{and press \kname{ENTER}}
\Do{Press}{\kname{SLASH}}{type \type{c}, press \kname{ENTER}, type
  \type{g6} and press \kname{ENTER}}
\Do{Press}{\kname{SLASH}}{type \type{c}, type \type{g5..g6} and press
  \kname{ENTER}, type \type{g7} and press \kname{ENTER}}
\Do{Press}{\kname{SLASH}}{type \type{c}, type \type{g5} and press
  \kname{ENTER}, type \type{g9} and press \kname{ENTER}}
\Do{Press}{\kname{SLASH}}{type \type{r}, \type{f}, \type{p}, type \type{1}
  and press \kname{ENTER}, type \type{g5..g9} and press \kname{ENTER}}
\end{Keystrokes}
The formula begins with a plus sign so that Initech 123 reads it as a
value.  Every copy divides by F9, and the shares appear as percentages
(Figure~\ref{fig:tl12}).
\FlairShot[fig:tl12]{tl12}{Each supply's share of the year}

\section{Lesson 13: Naming a Range}
\ix{range names}
A \term{range name} stands for a range in formulas and prompts.  Name the
grand total TOTAL, and use the name in the share formulas.

\begin{Keystrokes}
\Do{Press}{\fkey{Goto}{F5}}{and type \type{f9}, then press \kname{ENTER}}
\Do{Press}{\kname{SLASH}}{and type \type{r}, \type{n}, \type{c} to select
  Range Name Create}
\Do{Type}{\type{total}}{the name, and press \kname{ENTER}}
\Do{Press}{\kname{ENTER}}{to name the current cell, F9}
\Do{Press}{\fkey{Goto}{F5}}{and type \type{g5}, then press \kname{ENTER}}
\Do{Type}{\type{+f5/\$total}}{and press \kname{ENTER}}
\end{Keystrokes}
Copy G5 to G6, and G5..G6 to G7, as in Lesson~12.  A dollar sign before a
range name makes it an absolute reference.  The control panel shows the
formula with the cell the name stands for (Figure~\ref{fig:tl13}).
\FlairShot[fig:tl13]{tl13}{The share formula written with the name TOTAL}

\section{Lesson 14: Recalculating by Hand}
\ix{recalculation}
Initech 123 recalculates after every entry.  In a large worksheet you may
prefer to recalculate only when you ask.

\begin{Keystrokes}
\Do{Press}{\kname{SLASH}}{and type \type{w}, \type{g}, \type{r}, \type{m}
  to select Worksheet Global Recalculation Manual}
\Do{Press}{\fkey{Goto}{F5}}{and type \type{b5}, then press \kname{ENTER}}
\Do{Type}{\type{170}}{and press \kname{ENTER}}
\end{Keystrokes}
B5 changes, but the totals do not, and the \ind{CALC} indicator appears on
the status line (Figure~\ref{fig:tl14a}).
\FlairShot[fig:tl14a]{tl14a}{Manual recalculation: the totals wait}
\begin{Keystrokes}
\Do{Press}{\fkey{Calc}{F9}}{to recalculate}
\end{Keystrokes}
\FlairShot[fig:tl14b]{tl14b}{After CALC (F9)}
\begin{Keystrokes}
\Do{Press}{\kname{SLASH}}{and type \type{w}, \type{g}, \type{r}, \type{a}
  to return to automatic recalculation}
\end{Keystrokes}

\section{Lesson 15: Testing a Condition}
\ix{"@IF}
\fn{@IF} chooses between two results: \texttt{@IF(B9>380,"Yes","No")} gives
\texttt{Yes} if B9 is more than 380, and \texttt{No} otherwise.

\begin{Keystrokes}
\Do{Press}{\fkey{Goto}{F5}}{and type \type{a11}, then press \kname{ENTER}}
\Do{Type}{\type{Over 380?}}{and press \kname{RIGHT}}
\Do{Type}{\type{@if(b9>380,"Yes","No")}}{and press \kname{ENTER}}
\Do{Press}{\kname{SLASH}}{type \type{c}, press \kname{ENTER}, type
  \type{c11} and press \kname{ENTER}}
\Do{Press}{\kname{SLASH}}{type \type{c}, type \type{b11..c11} and press
  \kname{ENTER}, type \type{d11} and press \kname{ENTER}}
\end{Keystrokes}
\FlairShot[fig:tl15]{tl15}{The quarters over 380}
Save the worksheet as REVIEW: press \kname{SLASH}, type \type{f} and
\type{s}, type \type{REVIEW} and press \kname{ENTER}.
